\section{生成気体は物質収支で区別する}\label{theory:inorganic:gas_generation}
酸を加えれば全て同じ気体が出るわけではない。脱水、酸塩基、酸化還元の三種類を分ける。
\[
 \mathrm{HCOOH}\longrightarrow\mathrm{CO}+\mathrm{H_2O}.
\]
ギ酸への濃硫酸の作用による脱水では一酸化炭素を生じ、二酸化炭素へ書き換えない\cite{analysis49_co_ref}。
NaClと濃硫酸の塩化水素、Znと希硫酸の水素、MnO2と濃塩酸の塩素、Cuと濃硝酸の二酸化窒素を区別する。塩素生成は
\[
 \mathrm{MnO_2}+4\mathrm{HCl}\longrightarrow
 \mathrm{MnCl_2}+\mathrm{Cl_2}+2\mathrm{H_2O}
\]
でMnの還元とClの酸化が同時に起こる\cite{inorganic_halogen_ref}。
濃硝酸と希硝酸、濃硫酸と希硫酸では酸化性や代表生成物が違うので、原典の濃度条件を保つ。

\section{カルシウム塩の溶解・水和・潮解}\label{theory:inorganic:calcium}
炭酸カルシウムは通常純水へ難溶だが、CO2を含む水では
\[
 \mathrm{CaCO_3}+\mathrm{CO_2}+\mathrm{H_2O}
 \rightleftharpoons\mathrm{Ca^{2+}}+2\mathrm{HCO_3^-}
\]
という平衡で溶解が促進される。酸を加えた場合は炭酸塩からCO2が放出される。
CO2の溶存、pH、溶解・析出は地下水の水質と石灰岩地形の説明につながる\cite{inorganic_waterchem}。
骨・歯の無機成分にはCaを含む水酸アパタイトがある。半水石膏は加水して二水石膏となり、結晶成長によって硬化する\cite{analysis49_gypsum_ref}。
塩化カルシウムの吸湿・潮解は空気中の水分を取り込む性質であり、石膏の硬化とは別の現象である。
